\section{Legal Implications of Resolution Choice}
\label{sec:legal}

The choice of geographic resolution carries legal implications that practitioners must consider.

\paragraph{County-preservation requirements.}
Option~B (tract$\to$county) uses county adjacency as a computational coarsening tool,
but this should not be conflated with legal county-preservation requirements.
Many state constitutions explicitly require minimizing county splits:
California Article~XXI~\S2(d), Colorado Article~V~\S47, Florida Article~III~\S21.
The multi-scale algorithm treats counties as computational units during coarse moves
but does not enforce county-preservation as a hard constraint.
Practitioners in county-preservation states must verify compliance separately using
\texttt{redist analyze --types splits} after plan generation.

\paragraph{VRA implications of resolution choice.}
The choice of block group (BG) vs.\ census tract resolution can affect minority
population percentages in near-threshold majority-minority districts.
A district that barely satisfies a 50\% minority VAP threshold at the tract level
may fall below threshold when redrawn at the BG level due to finer boundary placement.
Practitioners deploying Option~A (BG$\to$tract) should recompute minority VAP
percentages at the BG level and verify VRA compliance using
\texttt{redist analyze --types bloc-voting} with BG-level assignments.

\paragraph{Resolution and the uniform distribution.}
Different resolution choices produce plans over different state spaces
(block-group partitions vs.\ tract partitions are distinct mathematical objects).
An ensemble produced at tract resolution and an ensemble produced at BG resolution
are not comparable distributions --- they answer different questions.
Expert witnesses should clearly specify the resolution level in any testimony
about ensemble distributions.
